\documentclass[11pt,a4paper]{article}

\usepackage[a4paper,margin=2.2cm]{geometry}
\usepackage[T1]{fontenc}
\usepackage[utf8]{inputenc}
\usepackage{lmodern}
\usepackage{microtype}
\usepackage{enumitem}
\usepackage{array}
\usepackage{booktabs}
\usepackage{tabularx}
\usepackage{longtable}
\usepackage{titlesec}
\usepackage[table]{xcolor}
\usepackage{graphicx}
\usepackage{fancyhdr}
\usepackage{hyperref}
\usepackage{multirow}
\usepackage{ragged2e}

\definecolor{kmutnbBlue}{HTML}{1F3A8A}
\definecolor{accent}{HTML}{0F6FC5}
\definecolor{lightGray}{HTML}{F2F4F7}
\definecolor{examrow}{HTML}{FEF3C7}
\definecolor{midrow}{HTML}{FEE2E2}

\hypersetup{
    colorlinks=true,
    linkcolor=accent,
    urlcolor=accent,
    citecolor=accent,
    pdfauthor={Asst. Prof. Dr. Ruslee Sutthaweekul},
    pdftitle={Course Syllabus -- 010153002 Computer Programming},
}

\titleformat{\section}
  {\normalfont\large\bfseries\color{kmutnbBlue}}
  {\thesection}{0.6em}{}
\titlespacing*{\section}{0pt}{1.2\baselineskip}{0.4\baselineskip}

\setlist[itemize]{leftmargin=1.4em,itemsep=2pt,topsep=2pt}
\setlist[enumerate]{leftmargin=1.6em,itemsep=2pt,topsep=2pt}

\pagestyle{fancy}
\fancyhf{}
\fancyhead[L]{\small\textit{010153002 Computer Programming}}
\fancyhead[R]{\small Semester 1/2026}
\fancyfoot[C]{\small\thepage}
\renewcommand{\headrulewidth}{0.3pt}

\newcommand{\fieldlabel}[1]{\textbf{#1:}}

\begin{document}

%% ─── Title block ────────────────────────────────────────────────────────────
\begin{center}
  {\Large\bfseries\color{kmutnbBlue}
    Electrical and Computer Engineering Course Syllabus}\\[4pt]
  {\large\bfseries 010153002 \quad Computer Programming}\\[2pt]
  {\normalsize Semester 1 / Academic Year 2026}\\[2pt]
  {\small Department of Electrical and Computer Engineering (ECE)}\\
  {\small Faculty of Engineering,
    King Mongkut's University of Technology North Bangkok}\\[4pt]
  {\small Instructor: Asst.\ Prof.\ Dr.\ Ruslee Sutthaweekul
    \quad\textbar\quad
    \href{mailto:ruslee.s@eng.kmutnb.ac.th}{ruslee.s@eng.kmutnb.ac.th}}
\end{center}

\noindent\rule{\linewidth}{0.4pt}\medskip

%% ─── 1. Course Information ──────────────────────────────────────────────────
\section*{1.\quad Course Information}
\begin{tabular}{@{}p{4.4cm}p{10.8cm}@{}}
  \fieldlabel{Course Number}   & 010153002 \\[2pt]
  \fieldlabel{Course Name}     & Computer Programming \\[2pt]
  \fieldlabel{Credits}         & 3 (2--2--5) \\[2pt]
  \fieldlabel{Contact Hours}   & 4 hours per week (2 lecture + 2 laboratory) \\[2pt]
  \fieldlabel{Schedule}        & Every Thursday, 13:00--17:00 \\[2pt]
  \fieldlabel{Semester Dates}  & 22 June 2026 -- 11 October 2026 \\[2pt]
  \fieldlabel{Course Type}     & Required \\[2pt]
  \fieldlabel{Prerequisites}   & None \\[2pt]
  \fieldlabel{Language}        & English (supplementary explanations in Thai available) \\[2pt]
  \fieldlabel{Textbook}        & Kernighan, B.~W.\ \& Ritchie, D.~M.\
                                  \textit{The C Programming Language},
                                  2nd Edition. Prentice Hall, 1988. \\
\end{tabular}

%% ─── 2. Course Description ──────────────────────────────────────────────────
\section*{2.\quad Course Description}
Basic operations of a computer, required components of a computer, how hardware
and software work together, data processing, fundamental of high level computer
programming, application design and development, problem solving by using
computer programming.

%% ─── 3. Course Learning Outcomes ────────────────────────────────────────────
\section*{3.\quad Course Learning Outcomes (CLOs)}
Upon successful completion of this course, students will be able to:
\begin{enumerate}
  \item Understand computer systems concepts, different computing environments,
        computer languages, and the system development life cycle.
  \item Write programs using basic programming constructs including identifiers,
        data types, variables, constants, input/output statements, logical and
        comparative operators, and control structures (selections and loops).
  \item Understand and apply expressions, precedence rules, associativity in
        evaluating expressions, type conversions, and various statement types.
  \item Develop and implement programs using arrays, pointers, and functions,
        including sorting algorithms, search algorithms, and string manipulations.
  \item Understand and utilize complex data structures such as structures, unions,
        and enumerated types.
  \item Recognize the importance of effective communication, ethical and
        professional responsibilities in engineering, and the impact of engineering
        solutions in various contexts.
\end{enumerate}

%% ─── 4. ABET Student Outcomes ───────────────────────────────────────────────
\section*{4.\quad ABET Student Outcomes Addressed}
\renewcommand{\arraystretch}{1.35}
\begin{tabularx}{\linewidth}{X >{\centering\arraybackslash}p{3.2cm}}
  \toprule
  \rowcolor{kmutnbBlue!90}
  \textcolor{white}{\textbf{ABET Student Outcome (SO)}} &
  \textcolor{white}{\textbf{Course Learning Outcome (CLO)}} \\
  \midrule
  1.\ An ability to identify, formulate, and solve complex engineering problems
      by applying principles of engineering, science, and mathematics.
  & 1--6 \\
  \rowcolor{lightGray}
  2.\ An ability to apply engineering design to produce solutions that meet
      specified needs with consideration of public health, safety, and welfare,
      as well as global, cultural, social, environmental, and economic factors.
  & 1--6 \\
  3.\ An ability to communicate effectively with a range of audiences.
  & 6 \\
  \rowcolor{lightGray}
  4.\ An ability to recognize ethical and professional responsibilities in
      engineering situations and make informed judgments, which must consider
      the impact of engineering solutions in global, economic, environmental,
      and societal contexts.
  & 6 \\
  5.\ An ability to function effectively on a team whose members together
      provide leadership, create a collaborative and inclusive environment,
      establish goals, plan tasks, and meet objectives.
  & --- \\
  \rowcolor{lightGray}
  6.\ An ability to develop and conduct appropriate experimentation, analyze
      and interpret data, and use engineering judgment to draw conclusions.
  & 1--6 \\
  7.\ An ability to acquire and apply new knowledge as needed, using
      appropriate learning strategies.
  & 1--6 \\
  \bottomrule
\end{tabularx}

%% ─── 5. CLO–Assessment Mapping ──────────────────────────────────────────────
\section*{5.\quad CLO Assessment Mapping}
\renewcommand{\arraystretch}{1.3}
\begin{tabularx}{\linewidth}{X >{\centering\arraybackslash}p{2.0cm} >{\centering\arraybackslash}p{4.4cm}}
  \toprule
  \rowcolor{kmutnbBlue!90}
  \textcolor{white}{\textbf{Course Learning Outcome (CLO)}} &
  \textcolor{white}{\textbf{SO}} &
  \textcolor{white}{\textbf{Assessment Tools}} \\
  \midrule
  1.\ Understand computer systems concepts, different computing environments,
      computer languages, and the system development life cycle.
  & 1, 7 & Exams, Quiz, Assignment, Project \\
  \rowcolor{lightGray}
  2.\ Write programs using basic programming constructs including identifiers,
      data types, variables, constants, input/output statements, logical and
      comparative operators, and control structures (selections and loops).
  & 1, 2, 7 & Exams, Quiz, Assignment, Project \\
  3.\ Understand and apply expressions, precedence rules, associativity in
      evaluating expressions, type conversions, and various statement types.
  & 1, 2, 7 & Exams, Quiz, Assignment, Project \\
  \rowcolor{lightGray}
  4.\ Develop and implement programs using arrays, pointers, and functions,
      including sorting algorithms, search algorithms, and string manipulations.
  & 1, 2, 7 & Exams, Quiz, Assignment, Project \\
  5.\ Understand and utilize complex data structures such as structures, unions,
      and enumerated types.
  & 1, 7 & Exams, Quiz, Assignment, Project \\
  \rowcolor{lightGray}
  6.\ Recognize the importance of effective communication, ethical and
      professional responsibilities in engineering, and the impact of engineering
      solutions in various contexts.
  & 1, 7 & Exams, Quiz, Assignment, Project \\
  \bottomrule
\end{tabularx}

%% ─── 6. Weekly Schedule ─────────────────────────────────────────────────────
\section*{6.\quad Weekly Schedule}

\noindent All classes meet on \textbf{Thursday, 13:00--17:00}.\quad
Lecture: 13:00--15:00\quad$|$\quad Laboratory: 15:00--17:00.

\vspace{6pt}
\renewcommand{\arraystretch}{1.25}
\setlength{\tabcolsep}{4pt}

\begin{longtable}{@{}
    >{\centering\bfseries}p{1.0cm}
    >{\centering}p{2.4cm}
    >{\RaggedRight}p{5.2cm}
    >{\RaggedRight}p{5.0cm}
  @{}}
  \toprule
  \rowcolor{kmutnbBlue!90}
  \textcolor{white}{\textbf{Week}} &
  \textcolor{white}{\textbf{Date}} &
  \textcolor{white}{\textbf{Lecture} \small(13:00--15:00)} &
  \textcolor{white}{\textbf{Lab} \small(15:00--17:00)} \\
  \midrule
  \endfirsthead
  \toprule
  \rowcolor{kmutnbBlue!90}
  \textcolor{white}{\textbf{Week}} &
  \textcolor{white}{\textbf{Date}} &
  \textcolor{white}{\textbf{Lecture} \small(13:00--15:00)} &
  \textcolor{white}{\textbf{Lab} \small(15:00--17:00)} \\
  \midrule
  \endhead
  \midrule
  \multicolumn{4}{r}{\small\itshape continued on next page}\\
  \endfoot
  \bottomrule
  \endlastfoot

  1 & Thu 25 Jun 2026
    & \textbf{Lec 1:} Introduction to Computer \& Programming\newline
      \small Computer components, hardware/software interaction, compilation pipeline,
             development environment.
    & \textbf{Lab 1:} Introduction to C\newline
      \small Compile and run \texttt{hello\_world.c}; explore compiler flags. \\

  \rowcolor{lightGray}
  2 & Thu 02 Jul 2026
    & \textbf{Lec 2:} Structure of a Program \& Data Types\newline
      \small Identifiers, variables, constants; integer, float, char;
             \texttt{printf}/\texttt{scanf}; type conversions.
    & \textbf{Lab 2:} Expressions \& Operators\newline
      \small Arithmetic operators, precedence rules, pre/post increment. \\

  3 & Thu 09 Jul 2026
    & \textbf{Lec 3:} Expressions \& Operators\newline
      \small Arithmetic, logical, and comparative operators; associativity;
             precedence; various statement types.
    & \textbf{Lab 3:} Expressions Practice\newline
      \small Type conversion exercises; operator precedence tasks. \\

  \rowcolor{lightGray}
  4 & Thu 16 Jul 2026
    & \textbf{Lec 4:} Selection -- Making Decisions\newline
      \small \texttt{if}, \texttt{if-else}, nested \texttt{if}, \texttt{switch};
             ternary operator.
    & \textbf{Lab 4:} Conditional Statements\newline
      \small Grade classifier, switch on user input. \\

  5 & Thu 23 Jul 2026
    & \textbf{Lec 5:} Repetition\newline
      \small \texttt{while}, \texttt{do-while}, \texttt{for}; \texttt{break},
             \texttt{continue}; nested loops.
    & \textbf{Lab 5:} Loop Control\newline
      \small Summation, factorial, pattern printing, input validation. \\

  \rowcolor{lightGray}
  6 & Thu 30 Jul 2026
    & \textbf{Lec 6:} Functions \& Macros\newline
      \small Declaration, definition, call, return values; scope; recursion;
             \texttt{\#define} macros.
    & \textbf{Lab 6:} Functions\newline
      \small Write and call helper functions; implement recursive factorial. \\

  7 & Thu 06 Aug 2026
    & \textbf{Lec 7:} Arrays\newline
      \small 1-D and 2-D arrays; traversal; passing arrays to functions;
             sorting and searching algorithms.
    & \textbf{Lab 7:} Arrays\newline
      \small Statistics on scores array; bubble sort; binary search. \\

  \rowcolor{lightGray}
  8 & Thu 13 Aug 2026
    & Part 1 Review \& Consolidation\newline
      \small Q\&A on Lectures 1--7; past exam walkthrough; common pitfalls.
    & \textbf{Lab 8:} Pointers (Part A)\newline
      \small Address-of \& dereference; pointer tracing; swap via pointers.
      \newline\small\textit{(Lec 8 content distributed as pre-reading.)} \\

  % ── Midterm exam period ────────────────────────────────────────────────────
  \rowcolor{midrow}
  \multicolumn{1}{>{\centering\bfseries\color{kmutnbBlue}}p{1.0cm}}{---}
    & \textbf{Thu 20 Aug 2026}
    & \multicolumn{2}{p{10.5cm}}{%
        \textbf{\color{kmutnbBlue}University Midterm Examination Period
        \quad (17--23 August 2026)}%
        \newline
        \small \textbf{No class this week.}\quad
        Midterm exam date \& venue announced by the Registrar.\quad
        Covers: Lectures 1--7, Labs 1--7.
      } \\

  9 & Thu 27 Aug 2026
    & \textbf{Lec 8:} Pointers\newline
      \small Address-of \& dereference; pointer arithmetic;
             \texttt{malloc}/\texttt{free}; dynamic memory.
    & \textbf{Lab 8:} Pointers (continued)\newline
      \small Dynamic array allocation; pointer to function. \\

  \rowcolor{lightGray}
  10 & Thu 03 Sep 2026
    & \textbf{Lec 9:} Strings\newline
      \small Null-terminated \texttt{char} arrays; \texttt{strlen}, \texttt{strcpy},
             \texttt{strcmp}, \texttt{strcat}; pointer traversal; string manipulation.
    & \textbf{Lab 9:} Strings\newline
      \small \texttt{my\_strlen}; reverse in place; palindrome checker;
             count vowels.
      \newline\small\textit{Project assignment announced.} \\

  11 & Thu 10 Sep 2026
    & \textbf{Lec 10:} Enumerated, Structure, and Union Types\newline
      \small \texttt{struct} definition; \texttt{.} vs \texttt{->};
             pass-by-pointer; \texttt{typedef}; \texttt{union}; \texttt{enum};
             memory padding.
    & \textbf{Lab 10:} Structures \& Unions\newline
      \small Student record with \texttt{->}; sort structs;
             tagged union area calculator. \\

  \rowcolor{lightGray}
  12 & Thu 17 Sep 2026
    & Project Work Session\newline
      \small Instructor and TA available for project Q\&A;
             individual code review on request.
    & Course Review \& Exam Preparation\newline
      \small Past exam questions; common bugs; tracing exercises. \\

  13 & Thu 24 Sep 2026
    & \multicolumn{2}{p{10.5cm}}{%
        Course Wrap-up \& Q\&A\newline
        \small Open Q\&A on all topics; exam strategy; feedback session.\quad
        \textbf{Last class.}\quad
        \textit{Project submission deadline: 17:00 today.}
      } \\

  % ── Final exam ────────────────────────────────────────────────────────────
  \rowcolor{examrow}
  \multicolumn{1}{>{\centering\bfseries\color{kmutnbBlue}}p{1.0cm}}{---}
    & \textbf{TBD}
    & \multicolumn{2}{p{10.5cm}}{%
        \textbf{\color{kmutnbBlue}Final Examination} \quad
        \textit{Date \& venue announced by the Registrar}%
        \newline
        \small Covers \textbf{Part 2 only}: Lectures 8--10 (Pointers, Strings,
               Structs/Unions) and Labs 8--10.\quad
               Closed-book; no computer; no AI tools.
      } \\

\end{longtable}

\noindent\textit{Note:} Teaching concludes on \textbf{24 September 2026} (Week~13).
The midterm exam period (17--23 Aug) and final exam dates are set by the Registrar.
The \textbf{final exam covers Part 2 only} (Lectures 8--10).

%% ─── 7. Assessment & Grading ────────────────────────────────────────────────
\section*{7.\quad Assessment and Grading}

\noindent
\begin{minipage}[t]{0.52\linewidth}
  \textbf{Assessment Breakdown}
  \begin{itemize}
    \item Lab Exercises, Quizzes \& Assignment \dotfill 40\%
    \item Midterm Examination          \dotfill 30\%
    \item Final Examination            \dotfill 30\%
  \end{itemize}
  \smallskip
  \noindent\small A minimum attendance of 80\,\% (11 of 13 sessions) is
  required to be eligible to sit the Final Examination.
\end{minipage}\hfill
\begin{minipage}[t]{0.44\linewidth}
  \textbf{Grading Scale}\par\smallskip
  \begin{tabular}{@{}l@{\hspace{1em}}l@{\hspace{2.5em}}l@{\hspace{1em}}l@{}}
    80--100 & A  & 60--64 & C  \\
    75--79  & B+ & 55--59 & D+ \\
    70--74  & B  & 50--54 & D  \\
    65--69  & C+ & 0--49  & F  \\
  \end{tabular}
\end{minipage}

%% ─── 8. Brief List of Topics ────────────────────────────────────────────────
\section*{8.\quad Topics Covered}
\begin{enumerate}[label=\alph*.]
  \item Introduction to Computer
  \item Introduction to Computer Language
  \item Structure of a Program
  \item Selection -- Making Decisions
  \item Repetition
  \item Arrays
  \item Pointers
  \item Functions
  \item Strings
  \item Enumerated, Structure, and Union Types
\end{enumerate}

%% ─── 9. Course Policies ─────────────────────────────────────────────────────
\section*{9.\quad Course Policies}

\noindent\textbf{Examinations.}
Both examinations are \textbf{closed-book and paper-based} (no computer, no AI tools).
Students must bring their student ID card.
Questions include multiple-choice, code tracing, bug fixing, and short
programming tasks written by hand.
No make-up exams except with documented medical or official university reasons
approved \emph{before} the exam date wherever possible.

\medskip
\noindent\textbf{Academic Integrity.}
Students are expected to submit their own work.
The use of AI tools to generate code submitted as individual work is not permitted.
Plagiarism and academic dishonesty are subject to university disciplinary
procedures which may include a grade of F for the course.

\medskip
\noindent\textbf{Lab Work.}
Lab sheets must be submitted by the end of the lab session.
Late submissions are not accepted without a valid excuse approved in advance.

\medskip
\noindent\textbf{Project.}
Individual assignment; collaboration on code is not permitted.
Submitted code must compile and run on GCC without modification.
Include a short written report (one page) describing your design choices.

\medskip
\noindent\textbf{Communication.}
Primary channel is the course LMS (announced in Week 1).
Email responses within two working days.

%% ─── 10. Course Materials ───────────────────────────────────────────────────
\section*{10.\quad Course Materials}
\begin{itemize}
  \item Lecture slides: \href{https://ruslylove.github.io/compro}{ruslylove.github.io/compro}
  \item Lab sheets: available on the course LMS before each class
  \item Compiler: GCC (Linux/macOS) or MinGW-GCC (Windows); any IDE is acceptable
  \item \textit{Reference:} Kernighan \& Ritchie,
        \textit{The C Programming Language}, 2nd ed.\ (Prentice Hall, 1988)
\end{itemize}

%% ─── Footer ─────────────────────────────────────────────────────────────────
\vfill
\noindent\rule{\linewidth}{0.4pt}
\begin{center}
  \small\textit{This syllabus is subject to minor adjustments at the
  instructor's discretion to accommodate the academic calendar and class
  progress. Students are responsible for staying informed of any changes
  announced in class or via the course LMS.}
\end{center}

\end{document}
